\section{Equilibrium Characterization}\label{sec:equilibrium}

\subsection{Product-market payoffs}

When both firms survive, the Cournot equilibrium is
\begin{equation}
x_D=p_D=\frac{1}{2+\gamma},
\qquad
\pi_D=\frac{1}{(2+\gamma)^2}.
\label{eq:duopoly}
\end{equation}
When only one firm survives,
\begin{equation}
x_M=p_M=\frac12,
\qquad
\pi_M=\frac14.
\label{eq:monopoly}
\end{equation}
Define the disaster-state scarcity-rent wedge
\begin{equation}
\Delta(\gamma)\equiv \pi_M-\pi_D
=\frac{\gamma(\gamma+4)}{4(\gamma+2)^2}.
\label{eq:delta}
\end{equation}
For $\gamma>0$, $\Delta>0$, and
\begin{equation}
\Delta'(\gamma)=\frac{2}{(\gamma+2)^3}>0.
\label{eq:delta-derivative}
\end{equation}
Thus closer substitutes increase the profit gain from being the sole surviving producer relative to competing in the disaster state.

\subsection{Backup choice}

Fix primary locations and public policies. Let $s_i$ and $s_j$ be the firms' primary failure probabilities and let $J$ be their joint primary failure probability. After integrating over disaster states, firm $i$'s expected profit can be written as
\begin{equation}
\Pi_i
=\pi_D-\pi_D s_i(1-r_i)
+\Delta s_j(1-r_j)
-\Delta J(1-r_i)(1-r_j)
-\frac{k}{2}r_i^2.
\label{eq:expected-profit}
\end{equation}
The interior first-order condition is
\begin{equation}
kr_i=\pi_D s_i+\Delta J(1-r_j).
\label{eq:backup-foc}
\end{equation}
The own second derivative is $-k<0$, so the best response is unique conditional on the rival's readiness.

\begin{proposition}[Strategic substitution in backup readiness]\label{prop:strategic-sub}
On the interior continuation branch,
\begin{equation}
\frac{\partial r_i^{BR}}{\partial r_j}=-\frac{\Delta J}{k}<0
\end{equation}
whenever $\gamma>0$ and $J>0$.
\end{proposition}

The economic force is specific to competition in the disruption state. A rival that remains productive after a common disruption eliminates the monopoly component of the firm's contingency payoff and therefore lowers its incentive to maintain backup readiness. Proposition~\ref{prop:strategic-sub} is a mechanism lemma, not a novelty claim.

For symmetric public risk $q$, the two location configurations deliver especially transparent expressions. If primary plants are co-located,
\begin{equation}
r^C(q)=\frac{q\pi_M}{k+q\Delta}.
\label{eq:rC}
\end{equation}
If primary plants are geographically dispersed,
\begin{equation}
r^D(q)=\frac{q(\pi_D+q\Delta)}{k+q^2\Delta}.
\label{eq:rD}
\end{equation}
Both are increasing in $q$, hence decreasing in public protection. Both are also decreasing in $\gamma$:
\begin{equation}
\frac{\partial r^C}{\partial \gamma}<0,
\qquad
\frac{\partial r^D}{\partial \gamma}<0.
\label{eq:gamma-r}
\end{equation}
The latter comparative static follows because a larger scarcity-rent wedge makes the value of backup more dependent on the rival's failure rather than on preserving competitive production in every state.

\subsection{Primary location}

Let $\Pi^{XY}_i$ denote firm $i$'s equilibrium continuation profit when firm $1$ locates in $X\in\{A,B\}$ and firm $2$ in $Y\in\{A,B\}$. For firm 1, define
\begin{align}
d_A&\equiv \Pi^{AA}_1-\Pi^{BA}_1,\\
d_B&\equiv \Pi^{AB}_1-\Pi^{BB}_1.
\end{align}
If the rival chooses $A$ with probability $p$, the deterministic payoff difference between choosing $A$ and $B$ is
\begin{equation}
D(p)=p d_A+(1-p)d_B.
\end{equation}
Because $\varepsilon_i$ is privately observed, the firm's cutoff rule implies the affine Bayesian best response
\begin{equation}
BR(p)=\frac12+\frac{D(p)}{2H}
=u_0+Bp,
\label{eq:location-br}
\end{equation}
where
\begin{equation}
u_0\equiv\frac{H+d_B}{2H},
\qquad
u_1\equiv\frac{H+d_A}{2H},
\qquad
B\equiv\frac{d_A-d_B}{2H}.
\label{eq:location-endpoints}
\end{equation}
Here $u_0=BR(0)$ and $u_1=BR(1)$. Whenever $u_0,u_1\in(0,1)$, the affine expression is valid globally on $p\in[0,1]$: no probability clipping occurs at either endpoint or anywhere between them. If in addition $|B|<1$, $BR$ is a contraction and the simultaneous private-information location game has a unique interior probability fixed point. Solving $p=BR(p)$ gives
\begin{equation}
p=\frac{H+d_B}{2H-d_A+d_B}.
\label{eq:location-prob}
\end{equation}
At the canonical witness, the exact whole-domain certificate verifies $0<u_0,u_1<1$ and $|B|<1$ for every $(a_A,a_B)\in[0,\bar a]^2$. Thus equation~\eqref{eq:location-prob} is the unique location continuation after every admissible first-stage deviation, not merely a local interior solution. At symmetric public policies $a_A=a_B$, symmetry implies $p=1/2$ and $D=0$. Off the symmetric policy path, equation~\eqref{eq:location-prob} captures the plant-attraction effect of a unilateral protection increase.
